\documentclass[11pt]{amsart}
\usepackage{amsmath,amssymb,amsthm,url,booktabs}
\newtheorem{theorem}{Theorem}\newtheorem{lemma}[theorem]{Lemma}
\theoremstyle{definition}\newtheorem{definition}[theorem]{Definition}
\theoremstyle{remark}\newtheorem{remark}[theorem]{Remark}
\newcommand{\MOAS}{\mathrm{MOAS}}
\title[Array-bic squares of order ten and of every large order]{Mutually orthogonal array-bic squares of order ten and of every large order}
\author{Yitzchak Shmalo}
\address{Einstein Institute of Mathematics, The Hebrew University of Jerusalem, Jerusalem, Israel}
\subjclass[2020]{05B15, 05B25, 05C70}
\keywords{array-bic square, regular antipodal plane, nearly generalised polygon, difference method, clique decomposition}
\begin{document}
\begin{abstract}
Adriaensen, Bamberg and Williford showed that $s$ mutually orthogonal array-bic squares of order $s$ are equivalent
to a regular antipodal plane of order $s$, found that three such squares do not exist for orders $5$, $6$ and $7$,
and asked whether three exist for some order greater than seven. They do: we construct three of order $10$, and
three of order $7n+3$ whenever $n$ is a product of prime powers each at least $5$. The squares of order $10$ extend
the unique triple of order $3$, and $10$ is the least order at which such an extension is possible. Using Keevash's
theorem on partite decompositions of typical graphs, we show that for each fixed $t$ there are $t$ mutually
orthogonal array-bic squares of every sufficiently large order, and that every such system embeds in one of every
sufficiently large order.
\end{abstract}
\maketitle

\section{Introduction}\label{sec:intro}
A nearly generalised $n$-gon of order $s$, in the sense of Payne~\cite{Payne}, is an $(s+1)$-regular bipartite graph
of girth $2n$ whose number of vertices exceeds the Moore bound by a small amount. Biggs and Ito~\cite{BiggsIto} proved
that such graphs do not exist for $s>1$ and $n>3$. For $n=3$ they are the incidence graphs of regular antipodal planes,
and Payne conjectured that a regular antipodal plane of order $s$ exists only for $s=2,3$. Adriaensen, Bamberg and
Williford~\cite{ABW} proved that none exists for $4\le s\le14$ and gave several reformulations of the problem. One of them uses array-bic squares (Definition~\ref{def:abs}): a
regular antipodal plane of order $s$ is equivalent to $s$ mutually orthogonal array-bic squares of order
$s$~\cite[Proposition~2.20]{ABW}. Writing $\MOAS(s,t)$ for a set of $t$ mutually orthogonal array-bic squares of
order $s$, they found by computer that $\MOAS(4,2)$ does not exist, that $\MOAS(s,2)$ exists and $\MOAS(s,3)$ does
not for $s=5,6,7$, and that the $\MOAS(3,3)$ is unique up to symmetry~\cite[Section~2.4]{ABW}. Among the questions
closing their paper is the following~\cite[Section~3, (Q3)]{ABW}.

\begin{quote}
Does there exist a $\MOAS(s,3)$ for $s>7$?
\end{quote}

We answer this question in the affirmative.

\begin{theorem}\label{thm:ten}
The three arrays of Table~\ref{tab:ten} form a $\MOAS(10,3)$, with positions $1$, $2$ and $3$. Their cells with
row and column in $\{1,2,3\}$ form the $\MOAS(3,3)$ of~\cite[Example~2.18]{ABW}.
\end{theorem}

\begin{theorem}\label{thm:family}
Let $n=1$ or let $n$ be a product of prime powers each at least $5$. Then a $\MOAS(7n+3,3)$ exists.
\end{theorem}

The hypothesis holds exactly when $2$ divides $n$ to the power $0$ or at least $3$ and $3$ divides $n$ to the power
$0$ or at least $2$; so every $n$ prime to $6$ satisfies it, and so do $8$, $9$, $16$ and $27$. The orders obtained
begin $10,38,52,59,66,80,94$. The squares of Theorem~\ref{thm:ten} extend the $\MOAS(3,3)$, and $10$ is the
least order at which this can happen, by the following bound (subsystems are defined in Section~\ref{sec:def}).

\begin{theorem}\label{thm:bound}
Let $t\ge2$. If a $\MOAS(s,t)$ contains a $\MOAS(u,t)$ as a subsystem and $u<s$, then $s\ge u+1+t(u-1)$.
\end{theorem}

For $u=t=3$ the bound is $s\ge10$. The existence of a $\MOAS(s,3)$ remains open at $s=8$ and $9$, and at every
order $s>10$ below the threshold of Theorem~\ref{thm:large} that is not covered by Theorem~\ref{thm:family}; by
Theorem~\ref{thm:bound} a $\MOAS(8,3)$ or a $\MOAS(9,3)$ cannot contain a $\MOAS(3,3)$ as a subsystem. For a
fixed number of squares the remaining orders are all covered once the order is large.

\begin{theorem}\label{thm:large}
Let $t\ge2$.
\begin{enumerate}
\item[(a)] There is $s_0(t)$ such that a $\MOAS(s,t)$ exists for every $s\ge s_0(t)$.
\item[(b)] For every $\MOAS(u,t)$ there is $s_1$ such that for every $s\ge s_1$ it is a subsystem of some $\MOAS(s,t)$,
with the same positions.
\end{enumerate}
\end{theorem}

Theorem~\ref{thm:large} uses Keevash's theorem on partite clique decompositions of typical balanced graphs~\cite{KeevashCD,KeevashII}
and gives no explicit value of $s_0(t)$. A regular antipodal plane of order $s$ corresponds to the case $t=s$, which
is far outside the range of Theorem~\ref{thm:large}. For $u=t$, Theorem~\ref{thm:bound} says that a
$\MOAS(t,t)$, that is a regular antipodal plane of order $t$, is a subsystem of a $\MOAS(s,t)$ with $s>t$ only if
$s\ge t^2+1$.

The squares of Theorem~\ref{thm:ten} are built by the difference method from four base cells: the three positions
and seven further points indexed by $\mathbb F_7$ are moved by an automorphism of order $3$, which permutes the
positions cyclically and multiplies $\mathbb F_7$ by $2$, and by the translations of $\mathbb F_7$
(Section~\ref{sec:ten}). For Theorem~\ref{thm:family} the group $\mathbb F_7$ is replaced by $\mathbb F_7\times H$,
where $H$ is an additive group of a product of finite fields, and the missing differences are supplied by the orbits of
cells whose five coordinates lie on a line of $\mathbb F_7\times H$ (Section~\ref{sec:family}).
Section~\ref{sec:bound} proves Theorem~\ref{thm:bound} and Section~\ref{sec:large} proves Theorem~\ref{thm:large}.

\section{Definitions}\label{sec:def}
By~\cite[Lemma~2.19]{ABW} the definitions depend on the point set only through its size, so we state them for an
arbitrary finite point set $P$ with $|P|=s$; in~\cite{ABW} the point set is $[s]=\{1,\dots,s\}$. The symbol $0$ is
not a point.

\begin{definition}[{\cite[Definition~2.16]{ABW}}]\label{def:abs}
Let $i\in P$. An \emph{array-bic square} on $P$ with \emph{position} $i$ is a map $A$ assigning to every pair
$(x,y)\in P^2$ with $x\ne y$ a symbol $A(x,y)\in(\{0\}\cup P)\setminus\{i\}$ such that
\begin{enumerate}
\item[(i)] $A(x,i)=0$ for every $x\ne i$;
\item[(ii)] for every $y\ne i$ the column $\{A(x,y):x\ne y\}$ contains every symbol of $P\setminus\{i\}$ exactly once;
\item[(iii)] for every $x$ the row $\{A(x,y):y\ne x\}$ contains every symbol of $P\setminus\{i,x\}$ exactly once and,
if $x\ne i$, also the symbol $0$.
\end{enumerate}
\end{definition}

\begin{definition}[{\cite[Definition~2.17]{ABW}}]\label{def:orth}
Array-bic squares $A_i$ and $A_k$ on $P$ with distinct positions $i$ and $k$ are \emph{orthogonal} if every pair
$(a,b)\in((\{0\}\cup P)\setminus\{i\})\times((\{0\}\cup P)\setminus\{k\})$ other than $(k,i)$ and the pairs $(c,c)$
with $c\in(\{0\}\cup P)\setminus\{i,k\}$ occurs as $(A_i(x,y),A_k(x,y))$ for exactly one pair $x\ne y$, and the
excluded pairs do not occur. A $\MOAS(s,t)$ is a set of $t$ pairwise orthogonal array-bic squares on a set of $s$
points; their positions are then distinct.
\end{definition}

A cell of an array-bic square is a pair $(x,y)$ of distinct points; a square has $s(s-1)$ cells. We record a cell of
a $\MOAS(s,t)$ with positions $i_1,\dots,i_t$ as the tuple $(x,y;A_{i_1}(x,y),\dots,A_{i_t}(x,y))$, and call its
entries its \emph{coordinates}: the row, the column and the outputs.

\begin{definition}\label{def:sub}
Let $\{A_i\}$ be a $\MOAS(s,t)$ on $P$, and let $U\subseteq P$ contain all positions. If $A_i(x,y)\in\{0\}\cup U$
for all distinct $x,y\in U$ and all $i$, and the restrictions of the $A_i$ to the cells in $U$ form a $\MOAS(|U|,t)$ on
$U$, we call these restrictions a \emph{subsystem} on $U$.
\end{definition}

All the conditions in Definitions~\ref{def:abs} and~\ref{def:orth} are statements about how often a given pair of values
occurs in a given pair of coordinates. We use them in that form. Let a finite family of tuples
$(x,y;a_1,\dots,a_t)$ with $x,y\in P$ and $a_1,\dots,a_t\in\{0\}\cup P$ be given, together with distinct positions
$\iota_1,\dots,\iota_t\in P$. For coordinates $\alpha,\beta$ and values $p,q$ let $N_{\alpha\beta}(p,q)$ be the number
of tuples with $\alpha$-coordinate $p$ and $\beta$-coordinate $q$. Write $R$ and $C$ for the row and column
coordinates, and $j$ for the output in position $\iota_j$. The family is the set of cells of a $\MOAS(s,t)$ with
these positions if and only if
\begin{equation}\label{eq:cell}
N_{RC}(x,y)=1 \text{ if } x\ne y,\qquad N_{RC}(x,x)=0,
\end{equation}
and, for all $j\ne k$,
\begin{align}
N_{Rj}(x,v)&=\begin{cases}1,& \text{if } v\in P\setminus\{\iota_j,x\}, \text{ or } v=0 \text{ and } x\ne\iota_j,\\
0,&\text{otherwise};\end{cases}\label{eq:row}\\
N_{Cj}(y,v)&=\begin{cases}1,& \text{if } y\ne \iota_j \text{ and } v\in P\setminus\{\iota_j\},\\
s-1,& \text{if } y=\iota_j \text{ and } v=0,\\ 0,&\text{otherwise};\end{cases}\label{eq:col}\\
N_{jk}(a,b)&=\begin{cases}1,& \text{if $(a,b)$ is a pair allowed in Definition~\ref{def:orth}},\\
0,&\text{otherwise}.\end{cases}\label{eq:orth}
\end{align}
Here $x$ and $y$ range over $P$ and $v$ over $\{0\}\cup P$. Indeed~\eqref{eq:cell} says that the tuples are indexed by
the cells, one tuple for each cell, and so define $t$ arrays $A_{\iota_j}(x,y)=a_j$. Then~\eqref{eq:row} says that row $x$ of $A_{\iota_j}$ consists of the symbols required by (iii), which are $s-1$ in
number, so it also forces every symbol to lie in $(\{0\}\cup P)\setminus\{\iota_j\}$;~\eqref{eq:col} is (i) and (ii),
since column $\iota_j$ has $s-1$ cells; and~\eqref{eq:orth} is Definition~\ref{def:orth}.

\section{Development over an abelian group}\label{sec:dev}
Let $G$ be a finite abelian group, written additively, and let $I=\{\iota_1,\iota_2,\iota_3\}$ be a set of three
points disjoint from $G$. The point set is $P=I\cup G$, so $s=|G|+3$, and the three squares have positions
$\iota_1,\iota_2,\iota_3$. Call the elements of $F=I\cup\{0\}$ \emph{fixed} and the elements of $G$ \emph{moving}. A
\emph{record} is a tuple $r=(r_R,r_C;r_1,r_2,r_3)$ with entries in $F\cup G$, and for $g\in G$ the translate
$r+g$ adds $g$ to every moving entry and keeps every fixed entry.

\begin{lemma}\label{lem:dev}
Let $B$ be a finite list of records with the following properties.
\begin{enumerate}
\item[(F)] Every record of $B$ has at most one fixed entry, except that for each $j\in\{1,2,3\}$ there is exactly one
record with $r_C=\iota_j$ and $r_j=0$, whose other three entries are moving. The fixed entries of the records with
exactly one fixed entry are, each exactly once: $r_R=\iota_x$ for $x=1,2,3$, and $r_j=\iota_k$ for the six pairs $j\ne k$.
\item[(D)] For every pair of coordinates $\alpha\ne\beta$ and every $\delta\in G$, the number of records of $B$ with
$r_\alpha,r_\beta\in G$ and $r_\beta-r_\alpha=\delta$ is $1$ if $\delta\ne0$; if $\delta=0$ it is $1$ when
$\{\alpha,\beta\}=\{C,j\}$ for some $j$, and $0$ otherwise.
\end{enumerate}
Let $T$ be the six tuples of the $\MOAS(3,3)$ of~\cite[Example~2.18]{ABW} on the points $\iota_1,\iota_2,\iota_3$ with
positions $\iota_1,\iota_2,\iota_3$. Then $T$ together with all translates $r+g$, $r\in B$, $g\in G$, is a
$\MOAS(|G|+3,3)$ on $I\cup G$, and $T$ is a subsystem on $I$.
\end{lemma}

\begin{proof}
First the count of tuples. By (D) for the pair $(R,C)$, the records with moving row and column number $|G|-1$; by (F)
there are three with fixed row and three with fixed column. So $|B|=|G|+5$ and there are
$6+|G|(|G|+5)=(|G|+2)(|G|+3)=s(s-1)$ tuples. We verify~\eqref{eq:cell}--\eqref{eq:orth}. Fix coordinates
$\alpha\ne\beta$ and values $p,q$.

\emph{Both values moving.} A translate $r+g$ has $\alpha$-coordinate $p$ and $\beta$-coordinate $q$ exactly when
$r_\alpha,r_\beta\in G$, $g=p-r_\alpha$ and $r_\beta-r_\alpha=q-p$. So $N_{\alpha\beta}(p,q)$ is the number in (D) for
$\delta=q-p$, and the tuples of $T$ do not contribute. The required values are: for $(R,C)$, $1$ if $q\ne p$ and $0$ if
$q=p$; for $(R,j)$, the same, since $q\in P\setminus\{\iota_j,p\}$ means $q\ne p$; for $(C,j)$, always $1$, since a
moving column is not $\iota_j$; for $(j,k)$, $1$ if $p\ne q$ and $0$ if $p=q$, since moving values are never
$\iota_j$ or $\iota_k$. These are the numbers in (D).

\emph{One value fixed.} If $p$ is fixed and $q$ is moving, each record with $r_\alpha=p$ and $r_\beta\in G$ has exactly
one translate with $\beta$-coordinate $q$, so $N_{\alpha\beta}(p,q)$ is the number of such records. By (F) the
records with a given fixed entry have all other entries moving except in the records with $r_C=\iota_j$, $r_j=0$. The
required values, all read from~\eqref{eq:row}--\eqref{eq:orth} with the fixed value in the coordinate indicated, are
\begin{itemize}
\item row $\iota_x$ against a moving column or output: $1$; this is the record with $r_R=\iota_x$;
\item column $\iota_y$ against a moving row: $1$; against a moving output $j$: $1$ if $j\ne y$ and $0$ if $j=y$; this is
the record with $r_C=\iota_y$, $r_y=0$, whose output $y$ is fixed;
\item output $j$ equal to $\iota_k$ against a moving row, column or output: $1$ if $k\ne j$ and $0$ if $k=j$; these are
the six records with $r_j=\iota_k$, $k\ne j$;
\item output $j$ equal to $0$ against a moving row or a moving output: $1$; against a moving column: $0$; this is the
record with $r_C=\iota_j$, $r_j=0$, whose column is fixed.
\end{itemize}
The case of $p$ moving and $q$ fixed is the same with the roles of $\alpha$ and $\beta$ exchanged.

\emph{Both values fixed.} Then $N_{\alpha\beta}(p,q)$ is the number of tuples of $T$ with these coordinates plus $|G|$
times the number of records of $B$ with $r_\alpha=p$ and $r_\beta=q$. By (F) the only records with two fixed entries
are those with $r_C=\iota_j$, $r_j=0$, one for each $j$; they contribute $|G|=s-3$ to $N_{Cj}(\iota_j,0)$ and nothing
elsewhere. The tuples of $T$ form a $\MOAS(3,3)$ on $I$, so they satisfy~\eqref{eq:cell}--\eqref{eq:orth} with $s=3$
and $P=I$, and they never use a moving value. Comparing with~\eqref{eq:cell}--\eqref{eq:orth} for $P=I\cup G$ restricted
to fixed values: the required counts are the same as for $s=3$, except $N_{Cj}(\iota_j,0)$, which is $s-1=2+|G|$
instead of $2$. The difference is supplied by the records just named. This proves the conditions for fixed values.

All cases are covered, so the family is a $\MOAS(s,3)$. The cells with row and column in $I$ are the tuples of $T$,
since no record has a fixed row and a fixed column, so $T$ is a subsystem on $I$.
\end{proof}

\section{Order ten}\label{sec:ten}
Take $G=\mathbb F_7$, and write $o_u$ for the point $u\in\mathbb F_7$, to distinguish it from the symbol $0$. Let
$\sigma$ be the permutation of $F\cup G$ that fixes $0$, maps $\iota_k\mapsto\iota_{k+1}$ (indices modulo $3$) and
$o_u\mapsto o_{2u}$, and for a record set
\[
\rho(r_R,r_C;r_1,r_2,r_3)=(\sigma r_R,\sigma r_C;\sigma r_3,\sigma r_1,\sigma r_2).
\]
Since $2^3=1$ in $\mathbb F_7$, $\rho^3$ is the identity. In the language of~\cite[Lemma~2.19]{ABW}, $\rho$ is the
relabelling of the points by $\sigma$, which carries the square in position $\iota_k$ to the square in position
$\iota_{k+1}$. Let $B$ consist of the four records
\[
(\iota_1,o_0;o_0,o_3,o_4),\quad(o_0,\iota_1;0,o_1,o_4),\quad(o_0,o_1;\iota_2,o_6,o_3),\quad(o_0,o_5;\iota_3,o_4,o_6)
\]
and their images under $\rho$ and $\rho^2$. The twelve records are listed in Table~\ref{tab:base}.

\begin{table}[ht]
\centering\small
\begin{tabular}{ccccc@{\qquad}ccccc}
\toprule
row & column & $1$ & $2$ & $3$ & row & column & $1$ & $2$ & $3$\\
\midrule
$\iota_1$ & $o_0$ & $o_0$ & $o_3$ & $o_4$ & $o_0$ & $o_1$ & $\iota_2$ & $o_6$ & $o_3$\\
$\iota_2$ & $o_0$ & $o_1$ & $o_0$ & $o_6$ & $o_0$ & $o_2$ & $o_6$ & $\iota_3$ & $o_5$\\
$\iota_3$ & $o_0$ & $o_5$ & $o_2$ & $o_0$ & $o_0$ & $o_4$ & $o_3$ & $o_5$ & $\iota_1$\\
$o_0$ & $\iota_1$ & $0$ & $o_1$ & $o_4$ & $o_0$ & $o_5$ & $\iota_3$ & $o_4$ & $o_6$\\
$o_0$ & $\iota_2$ & $o_1$ & $0$ & $o_2$ & $o_0$ & $o_3$ & $o_5$ & $\iota_1$ & $o_1$\\
$o_0$ & $\iota_3$ & $o_4$ & $o_2$ & $0$ & $o_0$ & $o_6$ & $o_2$ & $o_3$ & $\iota_2$\\
\bottomrule
\end{tabular}
\caption{The twelve base records for order $10$. In each half of the table the first and fourth lines are two of the
four records displayed above, and the two lines below each are its images under $\rho$ and $\rho^2$.}\label{tab:base}
\end{table}

\begin{proof}[Proof of Theorem~\ref{thm:ten}]
Condition (F) of Lemma~\ref{lem:dev} can be read from Table~\ref{tab:base}: the first record and its images give
$r_R=\iota_1,\iota_2,\iota_3$; the second and its images give $r_C=\iota_j$ with $r_j=0$; the third and fourth and
their images give $r_j=\iota_{j+1}$ and $r_j=\iota_{j-1}$; every other entry is moving. For (D), number the records
$1$--$6$ down the left block of Table~\ref{tab:base} and $7$--$12$ down the right block. Table~\ref{tab:diff} gives, for
each pair of coordinates $(\alpha,\beta)$ and each record whose $\alpha$- and $\beta$-coordinates are both moving, the
difference $r_\beta-r_\alpha\in\mathbb F_7$. In each of the three rows $(C,j)$ the entries are $0,1,\dots,6$, each once,
and in each of the other seven rows they are $1,2,\dots,6$, each once. This is condition (D).

\begin{table}[ht]
\centering\small
\begin{tabular}{lcccccccccccc}
\toprule
pair & $1$ & $2$ & $3$ & $4$ & $5$ & $6$ & $7$ & $8$ & $9$ & $10$ & $11$ & $12$\\
\midrule
$(R,C)$ & $\cdot$ & $\cdot$ & $\cdot$ & $\cdot$ & $\cdot$ & $\cdot$ & 1 & 2 & 4 & 5 & 3 & 6 \\
$(R,1)$ & $\cdot$ & $\cdot$ & $\cdot$ & $\cdot$ & 1 & 4 & $\cdot$ & 6 & 3 & $\cdot$ & 5 & 2 \\
$(R,2)$ & $\cdot$ & $\cdot$ & $\cdot$ & 1 & $\cdot$ & 2 & 6 & $\cdot$ & 5 & 4 & $\cdot$ & 3 \\
$(R,3)$ & $\cdot$ & $\cdot$ & $\cdot$ & 4 & 2 & $\cdot$ & 3 & 5 & $\cdot$ & 6 & 1 & $\cdot$ \\
$(C,1)$ & 0 & 1 & 5 & $\cdot$ & $\cdot$ & $\cdot$ & $\cdot$ & 4 & 6 & $\cdot$ & 2 & 3 \\
$(C,2)$ & 3 & 0 & 2 & $\cdot$ & $\cdot$ & $\cdot$ & 5 & $\cdot$ & 1 & 6 & $\cdot$ & 4 \\
$(C,3)$ & 4 & 6 & 0 & $\cdot$ & $\cdot$ & $\cdot$ & 2 & 3 & $\cdot$ & 1 & 5 & $\cdot$ \\
$(1,2)$ & 3 & 6 & 4 & $\cdot$ & $\cdot$ & 5 & $\cdot$ & $\cdot$ & 2 & $\cdot$ & $\cdot$ & 1 \\
$(1,3)$ & 4 & 5 & 2 & $\cdot$ & 1 & $\cdot$ & $\cdot$ & 6 & $\cdot$ & $\cdot$ & 3 & $\cdot$ \\
$(2,3)$ & 1 & 6 & 5 & 3 & $\cdot$ & $\cdot$ & 4 & $\cdot$ & $\cdot$ & 2 & $\cdot$ & $\cdot$ \\
\bottomrule
\end{tabular}
\caption{The differences $r_\beta-r_\alpha$ of the twelve base records; a dot marks a record in which one of the two
coordinates is fixed.}\label{tab:diff}
\end{table}

Lemma~\ref{lem:dev}
now shows that the tuples of~\cite[Example~2.18]{ABW} together with the $84$ translates form a $\MOAS(10,3)$ on
$I\cup\mathbb F_7$ with the $\MOAS(3,3)$ as a subsystem on $I$. Writing $\iota_k$ as $k$ and $o_u$ as $u+4$ gives the
arrays of Table~\ref{tab:ten}.
\end{proof}

\begin{table}[ht]
\centering\footnotesize
\setlength{\arraycolsep}{2.6pt}
\[
A_1=\begin{array}{cccccccccc}
\cdot&3&2&4&5&6&7&8&9&10\\
0&\cdot&3&5&6&7&8&9&10&4\\
0&2&\cdot&9&10&4&5&6&7&8\\
0&5&8&\cdot&2&10&9&7&3&6\\
0&6&9&7&\cdot&2&4&10&8&3\\
0&7&10&3&8&\cdot&2&5&4&9\\
0&8&4&10&3&9&\cdot&2&6&5\\
0&9&5&6&4&3&10&\cdot&2&7\\
0&10&6&8&7&5&3&4&\cdot&2\\
0&4&7&2&9&8&6&3&5&\cdot
\end{array}\qquad
A_2=\begin{array}{cccccccccc}
\cdot&0&3&7&8&9&10&4&5&6\\
3&\cdot&1&4&5&6&7&8&9&10\\
1&0&\cdot&6&7&8&9&10&4&5\\
5&0&6&\cdot&10&3&1&9&8&7\\
6&0&7&8&\cdot&4&3&1&10&9\\
7&0&8&10&9&\cdot&5&3&1&4\\
8&0&9&5&4&10&\cdot&6&3&1\\
9&0&10&1&6&5&4&\cdot&7&3\\
10&0&4&3&1&7&6&5&\cdot&8\\
4&0&5&9&3&1&8&7&6&\cdot
\end{array}
\]
\[
A_3=\begin{array}{cccccccccc}
\cdot&2&0&8&9&10&4&5&6&7\\
1&\cdot&0&10&4&5&6&7&8&9\\
2&1&\cdot&4&5&6&7&8&9&10\\
8&6&0&\cdot&7&9&5&1&10&2\\
9&7&0&2&\cdot&8&10&6&1&4\\
10&8&0&5&2&\cdot&9&4&7&1\\
4&9&0&1&6&2&\cdot&10&5&8\\
5&10&0&9&1&7&2&\cdot&4&6\\
6&4&0&7&10&1&8&2&\cdot&5\\
7&5&0&6&8&4&1&9&2&\cdot
\end{array}
\]
\caption{A $\MOAS(10,3)$ with positions $1,2,3$. Rows and columns are indexed $1,\dots,10$; the entry in row $x$ and
column $y$ of $A_i$ is $A_i(x,y)$.}\label{tab:ten}
\end{table}

\begin{remark}\label{rem:aut}
By construction the system is invariant under the group generated by the translations of $\mathbb F_7$ and the
relabelling $\sigma$, acting as in~\cite[Lemma~2.19]{ABW}. This group is the Frobenius group of order $21$; it permutes
the three squares cyclically, and its orbits on the $90$ cells are the six cells in $I$, which form two orbits of size
three, and four orbits of size $21$, one for each of the four records in Table~\ref{tab:base}.
\end{remark}

\section{An infinite family}\label{sec:family}
Let $n=q_1\cdots q_m$ with each $q_l$ a prime power, $q_l\ge5$, and let $H=\mathbb F_{q_1}\times\dots\times\mathbb F_{q_m}$
as an additive group, so $|H|=n$; for $n=1$ the group $H$ is trivial and the construction below is that of
Section~\ref{sec:ten}. In each field $\mathbb F_{q_l}$ choose distinct elements $c_{l,0}=0,c_{l,1},\dots,c_{l,4}$, which
is possible since $q_l\ge5$, and write $c_k=(c_{1,k},\dots,c_{m,k})\in H$, acting on $H$ by coordinatewise
multiplication. Let $G=\mathbb F_7\times H$ and for $k=0,\dots,4$ define $\varphi_k\colon G\to G$ by
$\varphi_k(a,h)=(ka,c_kh)$.

\begin{lemma}\label{lem:lines}
For $0\le k<l\le4$ the map $\varphi_l-\varphi_k$ is a bijection of $G$ that maps $G^*=\{(a,h):h\ne0\}$ onto itself and
$\mathbb F_7\times\{0\}$ onto itself.
\end{lemma}

\begin{proof}
$(\varphi_l-\varphi_k)(a,h)=((l-k)a,(c_l-c_k)h)$. Since $1\le l-k\le4$, multiplication by $l-k$ is a bijection of
$\mathbb F_7$, and in each factor $\mathbb F_{q_j}$ multiplication by $c_{j,l}-c_{j,k}\ne0$ is a bijection. The map
preserves the set of $(a,h)$ with $h=0$, hence also its complement.
\end{proof}

\begin{proof}[Proof of Theorem~\ref{thm:family}]
Identify $\mathbb F_7$ with $\mathbb F_7\times\{0\}\subseteq G$, so that the twelve records of Table~\ref{tab:base} are
records over $G$. For $d\in G^*$ let
\[
E_d=(\varphi_0(d),\varphi_1(d);\varphi_2(d),\varphi_3(d),\varphi_4(d)),
\]
whose entries are all moving, and let $B$ consist of the twelve records of Table~\ref{tab:base} and the records $E_d$,
$d\in G^*$. We check the hypotheses of Lemma~\ref{lem:dev}. Condition (F) holds because it holds for the twelve
records and the records $E_d$ have no fixed entries. For (D), take coordinates $\alpha\ne\beta$, the $k$-th and the
$l$-th entries of a record with $k<l$ in the order $R,C,1,2,3$, and $\delta\in G$. The twelve records have their moving
entries in $\mathbb F_7\times\{0\}$, so by Section~\ref{sec:ten} they contribute one record with difference $\delta$
for each $\delta\in(\mathbb F_7\times\{0\})\setminus\{0\}$, and, when $\{\alpha,\beta\}=\{C,j\}$, one for $\delta=0$;
they contribute nothing for $\delta\in G^*$. The record $E_d$ has difference $(\varphi_l-\varphi_k)(d)$, and by
Lemma~\ref{lem:lines} each $\delta\in G^*$ arises from exactly one $d\in G^*$, while no $\delta$ with zero second
component arises. Together this is (D). Lemma~\ref{lem:dev} gives a $\MOAS(|G|+3,3)=\MOAS(7n+3,3)$.
\end{proof}

When $n$ is prime to $6$ one may take $H=\mathbb Z/n\mathbb Z$ with $c_k=k$, since $1,\dots,4$ are then invertible
modulo $n$; the proof is the same.

\section{The extension bound}\label{sec:bound}
\begin{proof}[Proof of Theorem~\ref{thm:bound}]
Let $\{A_i\}$ be a $\MOAS(s,t)$ on $P$ with a subsystem on $U$, $|U|=u<s$, and let $\iota_1,\dots,\iota_t\in U$ be the
positions. Fix a point $x\in P\setminus U$. For each $j$, row $x$ of $A_{\iota_j}$ contains every symbol of
$U\setminus\{\iota_j\}$, since these lie in $P\setminus\{\iota_j,x\}$; these are $u-1$ symbols, in $u-1$ distinct
cells $(x,y)$.

No such cell has $y\in U$. If $y=\iota_j$ the entry is $0$. If $y\in U\setminus\{\iota_j\}$, the column $y$ of the
subsystem square restricted to $U$ already contains every symbol of $U\setminus\{\iota_j\}$ in its $u-1$ cells with row
in $U$, so by condition (ii) the cell $(x,y)$ cannot contain another one.

For $j\ne k$ no cell $(x,y)$ has $A_{\iota_j}(x,y)\in U\setminus\{\iota_j\}$ and $A_{\iota_k}(x,y)\in U\setminus\{\iota_k\}$.
Such a pair $(a,b)$ is either excluded by Definition~\ref{def:orth} (when $a=b$ or $(a,b)=(\iota_k,\iota_j)$) or
allowed, and every allowed pair with both entries in $\{0\}\cup U$ already occurs in a cell of the subsystem, because
the subsystem is itself a $\MOAS(u,t)$ on $U$; in either case it cannot occur in the cell $(x,y)$.

Hence the $t(u-1)$ cells found above are distinct cells $(x,y)$ with $y\in P\setminus(U\cup\{x\})$, a set of $s-u-1$
columns. So $s-u-1\ge t(u-1)$.
\end{proof}

\section{Every large order}\label{sec:large}
We encode a $\MOAS(s,t)$ by a clique decomposition of a partite graph and apply the theorem of Keevash on partite
decompositions. We state the special case we use. Let $\Gamma$ be a graph whose vertex set is partitioned into parts
$V_1,\dots,V_q$, with every edge joining two different parts. For distinct parts let
$d_{ab}=e(V_a,V_b)/(|V_a|\,|V_b|)$. The graph is \emph{$K_q$-balanced} if $e(V_a,V_b)$ is the same for all pairs of
parts and every vertex has the same number of neighbours in each part other than its own. It is
\emph{$(c,h)$-typical} if for every $h'\le h$, all distinct vertices $v_1,\dots,v_{h'}$ and every part $V_a$ containing
none of them, with $v_j\in V_{a_j}$,
\[
\Bigl|V_a\cap\bigcap_{j\le h'}N(v_j)\Bigr|=(1\pm h'c)\,|V_a|\prod_{j\le h'}d_{a_ja}.
\]
A \emph{partite $K_q$-decomposition} is a partition of the edges into copies of $K_q$ with one vertex in each part.

\begin{theorem}[{Keevash~\cite[Theorem~2.4]{KeevashCD},~\cite[Theorem~1.7]{KeevashII}, for graphs and $\mathcal H=K_q$}]\label{thm:keevash}
Let $q\ge2$, $h=2^{50q^3}$ and $\delta=2^{-10^3q^5}$. There are $c_0=c_0(q)>0$ and $N_0=N_0(q)$ such that the following
holds. Let $N>N_0$, let $d>2N^{-\delta/h^q}$ and $0<c<c_0d^{h^{30q}}$, and let $\Gamma$ be a $K_q$-balanced,
$(c,h^q)$-typical graph with parts $V_1,\dots,V_q$ of sizes between $N/h$ and $N$ and with $d_{ab}>d$ for all pairs of
parts. Then $\Gamma$ has a partite $K_q$-decomposition.
\end{theorem}

In~\cite{KeevashCD} the theorem is stated for an $r$-graph $\mathcal H$ on $[q]$ and an $\mathcal H$-blowup; for $r=2$ and $\mathcal H=K_q$ the
typicality and balance conditions there are the ones above. We use only that $h$, $\delta$ and $c_0$ depend on $q$
alone.

\begin{proof}[Proof of Theorem~\ref{thm:large}(a)]
Let $P=[s]$ and let $I\subseteq P$, $|I|=t$, be the set of positions. Take disjoint vertex sets: $X=P$ (rows),
$Y=P\setminus I$ (ordinary columns), and for $j\in I$ a set $S_j$, a copy of $P\setminus\{j\}$ (the positive
symbols of the square in position $j$); a vertex of $S_j$ carrying the label $a$ is written $a$ when the part is
clear. Let $\Gamma_0$ be the $(t+2)$-partite graph with the following edges, for labels $x,y,a,b$:
\begin{align*}
&x\in X\sim y\in Y \iff x\ne y; && x\in X\sim a\in S_j\iff x\ne a;\\
&y\in Y\sim a\in S_j \text{ always}; && a\in S_j\sim b\in S_k\iff a\ne b,\ (a,b)\ne(k,j).
\end{align*}
An ordinary cell $(x,y)$, $y\in Y$, with outputs $a_j\in S_j$ satisfying the conditions of Definitions~\ref{def:abs}
and~\ref{def:orth} is a $(t+2)$-clique of $\Gamma_0$; a cell $(x,i)$ in the column of position $i$ has output $0$ in
square $i$ and corresponds to a $t$-clique on $X$ and the parts $S_j$, $j\ne i$. Conversely we shall choose, for
every $i\in I$, $s-1$ vertex-disjoint such $t$-cliques, one through every $x\ne i$ and one through every vertex of each
$S_j$, $j\ne i$, with all these cliques edge-disjoint; and then a partite $K_{t+2}$-decomposition of the remaining
graph. Such a choice gives a $\MOAS(s,t)$: each full clique is read as the cell indexed by its $X$--$Y$ edge, which
determines the cell uniquely, and the clique chosen through $x$ in step $i$ is the cell $(x,i)$, with output $0$ in
square $i$. Every $X$--$Y$ edge lies in exactly one clique, so every ordinary cell receives exactly one tuple. Each
$X$--$S_j$ edge lies in exactly one clique, which is~\eqref{eq:row} for positive symbols; the symbol $0$ of square $j$
appears in row $x\ne j$ exactly in the cell $(x,j)$. Each $Y$--$S_j$ edge lies in exactly one clique, and column $i$ of
square $j\ne i$ meets every vertex of $S_j$ once; this is~\eqref{eq:col}. Each $S_j$--$S_k$ edge lies in exactly one
clique, which gives~\eqref{eq:orth} for pairs of positive symbols, and the pairs $(0,b)$ and $(a,0)$ occur exactly in
the columns $j$ and $k$, once for each positive $b$, respectively $a$; the pair $(0,0)$ does not occur because a cell
has at most one zero.

\emph{The columns of the positions.} Treat the $i\in I$ one at a time, and for a fixed $i$ adjoin the parts $S_j$,
$j\in I\setminus\{i\}$, one at a time to the list of rows $x\ne i$. At each stage we have $s-1$ partial tuples and must
match them to the $s-1$ vertices of the next part $S_j$, so that every new pair is an edge of $\Gamma_0$ not used by a
previously chosen column. Call a vertex $w$ of $S_j$ \emph{forbidden} for a vertex $v$ of another part if $vw$ is not
an edge of $\Gamma_0$ or was used by an earlier column. From the definition of $\Gamma_0$, a vertex $v$ has at most one
non-neighbour in each other part, and each of the at most $t-1$ earlier columns used at most one edge at $v$ into
$S_j$; so at most $t$ vertices of $S_j$ are forbidden for $v$. A partial tuple has at most $t-1$ entries, so at most
$t(t-1)$ vertices of $S_j$ are forbidden for one of its entries. Conversely, a vertex $w$ of $S_j$ is forbidden for at
most $t$ vertices of each of the at most $t-1$ parts already in the tuples, by the same count, and each vertex of
such a part is an entry of exactly one partial tuple; so $w$ is forbidden for an entry of at most $t(t-1)$ partial
tuples. Within the current column the new edges are distinct, because the partial tuples have distinct entries in each
part and are matched to distinct vertices of $S_j$. The bipartite graph of allowed pairs between the $s-1$ partial
tuples and $S_j$ therefore has minimum degree at least $s-1-t(t-1)\ge s-t^2\ge(s-1)/2$ for $s\ge2t^2$, and such a balanced bipartite graph has a perfect matching by Hall's
theorem (a set $S$ of size at most $(s-1)/2$ has at least $(s-1)/2\ge|S|$ neighbours, and a larger set meets the
neighbourhood of every vertex on the other side). This chooses all the $t$-cliques.

\emph{Balance of the remaining graph.} Let $\Gamma$ be $\Gamma_0$ with the edges of the chosen $t$-cliques removed. A
vertex $y\in Y$ has $s-1$ neighbours in every other part. A vertex $a\in S_j$ had $s-1$ neighbours in $X$, of which
the $t-1$ columns $i\ne j$ used one each; $s-t$ in $Y$, none used; and $s-2$ in $S_k$ (the label $a$ is excluded if
$a\ne k$, the label $j$ if $a=k$), of which the $t-2$ columns $i\notin\{j,k\}$ used one each. So $a$ has $s-t$
neighbours in every other part. A row $x$ has $s-t-[x\notin I]$ neighbours in $Y$; in $S_j$ it had $s-1$ if $x=j$ and
$s-2$ otherwise, and the columns $i\in I\setminus\{j,x\}$ used one each, leaving $s-t-[x\notin I]$ again. Thus every
vertex has the same number of neighbours in each other part, and summing over a part gives
$e(V_a,V_b)=(s-t)(s-1)$ for every pair of parts. So $\Gamma$ is $K_{t+2}$-balanced.

\emph{Typicality.} Every part has between $s-t$ and $s$ vertices, and every vertex of $\Gamma$ has at least $s-t-1$
neighbours in each other part, so at most $t+1$ non-neighbours there. The parts have sizes $s$, $s-t$ and $s-1$, so the
densities $d_{ab}=(s-t)(s-1)/(|V_a|\,|V_b|)$ are $1-1/s$, $1-t/s$, $1$ and $(s-t)/(s-1)$, all in $[1-t/s,1]$. Let
$h'\le h^q$, with $q=t+2$, and suppose $s\ge2h^q(t+1)$. The product of $h'$ densities lies in $[1-h't/s,1]$, and the
common neighbourhood of $h'$ vertices in a part $V_a$ has between $|V_a|-h'(t+1)$ and $|V_a|$ elements. Since
$h'(t+1)/s\le1/2$ and $|V_a|\ge s-t\ge s/2$, the ratio of the two sides of the typicality identity lies in $[1-2h'(t+1)/s,\,1+2h'(t+1)/s]$. So
$\Gamma$ is $(c,h^q)$-typical as soon as also $s\ge2(t+1)/c$.

Now apply Theorem~\ref{thm:keevash} with $q=t+2$, $N=s$ and $d=1/2$, and fix $c$ with $0<c<c_0(1/2)^{h^{30q}}$. For
$s$ large, $1/2>2s^{-\delta/h^q}$, the parts have sizes between $s/h$ and $s$, all densities exceed $1/2$, and $\Gamma$
is $(c,h^q)$-typical. So $\Gamma$ has a partite $K_{t+2}$-decomposition, and the $\MOAS(s,t)$ follows.
\end{proof}

\begin{proof}[Proof of Theorem~\ref{thm:large}(b)]
By~\cite[Lemma~2.19]{ABW} we may assume that the given $\MOAS(u,t)$ is a system on $U=[u]$ and that $I\subseteq U$. Construct $\Gamma_0$ for
$P=[s]$ as above. The cells of the given system with row and column in $U$ are cliques of $\Gamma_0$: full cliques for
ordinary columns $y\in U\setminus I$, and $t$-cliques for the columns in $I$. They are edge-disjoint, because two of
them sharing an edge would repeat a row entry, a column entry or an output pair of the given system. Reserve them. In
the step for the column $i\in I$ the rows in $U\setminus\{i\}$ are already matched, to the labels in
$U\setminus\{j\}$ of each $S_j$, by condition (ii) for the given system; match the $s-u$ rows of $P\setminus U$ to the
$s-u$ labels of $S_j$ in $P\setminus U$ by the same greedy procedure. The reserved cliques use only vertices with
labels in $U$ and the new edges join vertices with labels in $P\setminus U$, so the counts of forbidden vertices are
as before, and the bipartite graph of allowed pairs has minimum degree at least $s-u-t^2+1\ge(s-u)/2$ for $s$ large.

Let $\Gamma$ be $\Gamma_0$ with all chosen $t$-cliques removed. The computation of degrees above did not depend on which
$t$-cliques were chosen, so $\Gamma$ is $K_{t+2}$-balanced. Let $\Gamma'$ be $\Gamma$ with the edges of the reserved full
cliques removed. A full clique has one vertex in every part, so removing it lowers the degree of each of its vertices
into each other part by exactly one, and lowers $e(V_a,V_b)$ by exactly one for every pair of parts. A vertex lying in
$m$ reserved cliques therefore loses $m$ neighbours in every other part, and $\Gamma'$ is still $K_{t+2}$-balanced. A
vertex lies in at most $u-1$ reserved cliques: a row or an ordinary column in $U$ lies in at most $u-1$ cells with row
and column in $U$, a symbol occurs at most once in each column, and there are at most $u-1$ ordinary columns in $U$.
So every vertex of $\Gamma'$ has at most $t+u$ non-neighbours in each other part. For $s\ge2t$ every part has at least
$s/2$ vertices, so every density of $\Gamma'$ is at least $1-2(t+u)/s$, and the computation for typicality above, with
these two bounds in place of $t+1$ and $1-t/s$, shows that $\Gamma'$ is $(c,h^q)$-typical once $s\ge4h^q(t+u)$ and
$s\ge4(t+u)/c$. Theorem~\ref{thm:keevash} gives a partite
$K_{t+2}$-decomposition of $\Gamma'$, and together with the reserved cliques this is a $\MOAS(s,t)$ whose cells in $U$
are those of the given system, which is therefore a subsystem on $U$.
\end{proof}

\section*{Acknowledgements and funding}
The research presented in this paper was supported by the European Research Council (ERC) under the European Union's
Horizon Europe research and innovation programme (grant agreement No.~101041711), by the Simons Foundation as part of
the Collaboration on the Mathematical and Scientific Foundations of Deep Learning, by Heights Labs, by Convex Nexus
Capital, by the Israel Science Foundation (grant number 2258/19), and by the Israel Science Foundation (ISF Grant
4101/25).

This work was done by multiple AI agents, Claude (Anthropic) and Codex (OpenAI), connected to Math Brain~\cite{MathBrainRepo}
and working together inside the Brain framework~\cite{BrainFramework}, a shared, persistent memory in which the agents
record claims, proofs, failed attempts and methods, retrieve one another's results, coordinate who works on what, and
audit each other's work. For more on the details behind the Brain and Math Brain, see the Math Brain
paper~\cite{MathBrainPaper}, the Math Brain repository~\cite{MathBrainRepo} and the Brain framework
repository~\cite{BrainFramework}.

\section*{Declaration of competing interest}
The author holds equity in Heights Labs and in Convex Nexus Capital and has a management role at Convex Nexus Capital.
These entities are included in the funding declaration above.

\section*{Declaration of generative AI and AI-assisted technologies}
Array-bic squares, their orthogonality, Question (Q3) and the $\MOAS(3,3)$ are due to Adriaensen, Bamberg and
Williford~\cite{ABW}; the decomposition theorem used in Section~\ref{sec:large} is due to Keevash~\cite{KeevashCD,KeevashII}.
A Codex agent (OpenAI) found the four base records of Section~\ref{sec:ten}, the family of Section~\ref{sec:family}, the
bound of Theorem~\ref{thm:bound} and the argument of Section~\ref{sec:large}. Claude Opus~5.5 (Anthropic) wrote the
proof of Lemma~\ref{lem:dev} and the proofs of Theorems~\ref{thm:ten} and~\ref{thm:family} in the form given here, the
programs that check them and the text. The author directed the work and is responsible for the content of the paper
and for decisions concerning its dissemination.

\section*{Data availability}
A program that builds the squares of Theorem~\ref{thm:ten}, and those of Theorem~\ref{thm:family} when each $q_l$ is a
prime, $8$ or $9$, from the base records and checks Definitions~\ref{def:abs} and~\ref{def:orth} directly is available
as tag v1.0 (commit 8990388) of the repository\\
\url{https://github.com/yspennstate/array-bic-squares-order-ten}.

\clearpage
\appendix
\section{Notation}
\begin{center}\scriptsize
\begin{tabular}{l p{6.4cm} p{2.6cm}}
\toprule
Symbol & Meaning & First defined\\
\midrule
$\MOAS(s,t)$ & $t$ mutually orthogonal array-bic squares of order $s$ & Section~\ref{sec:intro}\\
$P$, $s$ & point set, $|P|$ & Section~\ref{sec:def}\\
$0$ & the symbol zero, not a point & Section~\ref{sec:def}\\
$A_i$, $i$ & array-bic square, its position & Definition~\ref{def:abs}\\
$R$, $C$, $j$ & row, column and output coordinates of a tuple & Section~\ref{sec:def}\\
$N_{\alpha\beta}(p,q)$ & number of tuples with coordinates $\alpha,\beta$ equal to $p,q$ & Section~\ref{sec:def}\\
$G$, $I$, $\iota_k$ & abelian group, set of positions, a position & Section~\ref{sec:dev}\\
$F$ & fixed values $I\cup\{0\}$ & Section~\ref{sec:dev}\\
$r$, $r+g$, $B$ & record, its translate, list of base records & Section~\ref{sec:dev}\\
$T$ & the six tuples of the $\MOAS(3,3)$ & Lemma~\ref{lem:dev}\\
$o_u$ & the point $u\in\mathbb F_7$ & Section~\ref{sec:ten}\\
$\sigma$, $\rho$ & relabelling of order $3$, its action on records & Section~\ref{sec:ten}\\
$n$, $q_l$, $H$ & $s=7n+3$, prime powers with product $n$, $\prod\mathbb F_{q_l}$ & Section~\ref{sec:family}\\
$c_k$, $\varphi_k$, $G^*$, $E_d$ & scalars in $H$, maps $(a,h)\mapsto(ka,c_kh)$, $\{(a,h):h\ne0\}$, line record & Section~\ref{sec:family}\\
$U$, $u$ & point set of a subsystem, its size & Definition~\ref{def:sub}\\
$t$ & number of squares & Section~\ref{sec:intro}\\
$X$, $Y$, $S_j$ & parts for rows, ordinary columns, positive symbols & Section~\ref{sec:large}\\
$\Gamma_0$, $\Gamma$, $\Gamma'$ & encoding graph, after the columns of the positions, after reserved cliques & Section~\ref{sec:large}\\
$V_a$, $K_q$, $N(v)$ & a part, the complete graph on $q$ vertices, neighbourhood of a vertex $v$ & Section~\ref{sec:large}\\
$w$ & a vertex of the part $S_j$ being matched & Section~\ref{sec:large}\\
$d_{ab}$, $h$, $\delta$, $c$, $c_0$ & density of a pair of parts; constants of Theorem~\ref{thm:keevash} & Section~\ref{sec:large}\\
\bottomrule
\end{tabular}
\end{center}

\begin{thebibliography}{99}
\bibitem{ABW} S.~Adriaensen, J.~Bamberg, J.~Williford, On configurations that are nearly generalised polygons, preprint (2026), arXiv:2611.17855.
\bibitem{BiggsIto} N.~L.~Biggs, T.~Ito, Graphs with even girth and small excess, Math. Proc. Cambridge Philos. Soc. 88 (1980), 1--10.
\bibitem{KeevashII} P.~Keevash, The existence of designs II, preprint (2018), arXiv:1802.05900.
\bibitem{KeevashCD} P.~Keevash, Coloured and directed designs, in: Building Bridges II, Bolyai Soc. Math. Stud. 28, Springer, Berlin, 2019, 279--315; theorem numbers as in arXiv:1807.05770.
\bibitem{Payne} S.~E.~Payne, On the non-existence of a class of configurations which are nearly generalized $n$-gons, J. Combin. Theory Ser. A 12 (1972), 268--282.
\bibitem{BrainFramework} Y.~Shmalo, \emph{Brain Framework}, public software repository, 2026, \url{https://github.com/yspennstate/brain-framework-public}.
\bibitem{MathBrainRepo} Y.~Shmalo, \emph{Math Brain: public snapshot, code and tests}, 2026, \url{https://github.com/yspennstate/math-brain-public}.
\bibitem{MathBrainPaper} Y.~Shmalo, \emph{Math Brain: shared, auditable memory for multi-agent mathematical research}, October 2026, doi:10.5281/zenodo.24150837, \url{https://github.com/yspennstate/math-brain-public/tree/main/papers/math-brain-2026-10-24}.
\end{thebibliography}
\end{document}
